%% ── Chapter 2: Background ────────────────────────────────────
\chapter{Background}
\label{ch:background}

\section{Generative Adversarial Networks}

A Generative Adversarial Network \citep{goodfellow2014} consists of two
neural networks trained in a minimax game:
\begin{equation}
  \min_G \max_D\; \E_{x \sim p_{\mathrm{data}}}[\log D(x)]
    + \E_{z \sim p_z}[\log(1 - D(G(z)))]
  \label{eq:gan}
\end{equation}
where $G: \mathcal{Z} \to \mathcal{X}$ is the generator mapping latent noise
$z \in \mathcal{Z}$ to data space, and $D: \mathcal{X} \to [0,1]$ is the
discriminator.

\subsection{Conditional GANs}
Extending~\eqref{eq:gan}, a conditional GAN \citep{mirza2014} incorporates
auxiliary information $\mathbf{c}$:
\begin{equation}
  \min_G \max_D\; \E_{x,\mathbf{c}}[\log D(x|\mathbf{c})]
    + \E_{z,\mathbf{c}}[\log(1 - D(G(z|\mathbf{c})|\mathbf{c}))]
\end{equation}
In \caesar{}, $\mathbf{c} = \mathbf{d}_t$ is the real-time defense embedding,
yielding the \tagan{} formulation.

\section{Deep Reinforcement Learning}

\subsection{Q-Learning and DQN}
In a Markov Decision Process $(\mathcal{S}, \mathcal{A}, \mathcal{R},
\mathcal{T}, \gamma)$, the Q-function satisfies the Bellman equation:
\begin{equation}
  Q^*(s,a) = \E\bigl[r + \gamma \max_{a'} Q^*(s',a') \mid s,a\bigr]
\end{equation}
DQN \citep{mnih2015} approximates $Q^*$ with a neural network $Q_\theta$,
minimising:
\begin{equation}
  \mathcal{L}(\theta) = \E_{(s,a,r,s') \sim \mathcal{D}}
    \Bigl[(r + \gamma \max_{a'} Q_{\hat\theta}(s',a') - Q_\theta(s,a))^2\Bigr]
\end{equation}
where $\mathcal{D}$ is a replay buffer and $\hat\theta$ are target network
parameters updated periodically.

\subsection{Dueling DQN}
The Dueling architecture \citep{wang2016} decomposes:
\begin{equation}
  Q(s,a;\theta,\alpha,\beta) = V(s;\theta,\beta)
    + A(s,a;\theta,\alpha) - \frac{1}{|\mathcal{A}|}\sum_{a'} A(s,a';\theta,\alpha)
\end{equation}
allowing the value stream $V(s)$ and advantage stream $A(s,a)$ to be learned
independently, which accelerates convergence in the \adpn{}.

\section{Denoising Diffusion Probabilistic Models}

A DDPM \citep{ho2020} defines a forward process:
\begin{equation}
  q(\mathbf{x}_t \mid \mathbf{x}_0) = \mathcal{N}
    \!\left(\mathbf{x}_t;\, \sqrt{\bar\alpha_t}\,\mathbf{x}_0,\,
    (1-\bar\alpha_t)\mathbf{I}\right)
  \label{eq:ddpm_forward}
\end{equation}
with $\bar\alpha_t = \prod_{s=1}^t (1 - \beta_s)$ and a cosine schedule
\citep{nichol2021}:
\begin{equation}
  \bar\alpha_t = \frac{f(t)}{f(0)},\quad
  f(t) = \cos^2\!\left(\frac{t/T + s}{1+s} \cdot \frac{\pi}{2}\right)
\end{equation}
The reverse process $p_\theta(\mathbf{x}_{t-1}|\mathbf{x}_t)$ is learned
by minimising:
\begin{equation}
  \mathcal{L}_{\mathrm{simple}} = \E_{t,\mathbf{x}_0,\boldsymbol\epsilon}
    \!\left[\|\boldsymbol\epsilon - \boldsymbol\epsilon_\theta(\mathbf{x}_t,t)\|^2\right]
\end{equation}
\mgae{} exploits this process to project attack samples onto the learned
benign traffic manifold (\cref{ch:mgae}).

\section{Graph Neural Networks and Knowledge Graphs}

The \tagraph{} maintains a directed graph $G = (V, E, w)$ where nodes represent
attack types and defense actions, and edge weights encode co-occurrence
statistics. PageRank \citep{page1999} identifies high-centrality attack
nodes:
\begin{equation}
  \mathrm{PR}(v) = \frac{1-d}{|V|} + d \sum_{u \in \mathcal{N}^-(v)}
    \frac{\mathrm{PR}(u)}{|\mathcal{N}^+(u)|}
\end{equation}
Markov-chain transition probabilities are estimated from edge weights to
enable attack sequence prediction (\cref{ch:tag}).
